\chapter{Tier 2: black-box experts and the contraction contract}

At tier 2 an expert exposes only its boundary response.  The host exchanges
wave variables between neighbouring sides and iterates.  Nothing is automatic
for a learned black box, so the guarantee is conditional: it holds when a
contraction factor has been \emph{certified} for the experts.  This chapter
states what follows from such a certificate (T8, T9a, T9b), and two examples
that show which hypothesis the certificate must be about.

\section{T8: wave variables and the Cayley transform}

\paragraph{In plain words.}
On a side, the effort $e$ (a temperature) and the flow $f$ (a flux into the
piece) are combined into an incoming wave $e+Zf$ and an outgoing wave $e-Zf$,
for an impedance $Z>0$.  The squared size of the outgoing wave is that of the
incoming one minus four times $Z$ times the power the piece absorbs.  So a
piece that never generates power does not amplify waves, and a piece that
absorbs a definite share of what it is given shrinks them by a definite factor.

\begin{lemma}[T8, the wave identity]\label{lem:T8-identity}
  \lean{Atlas.wave_identity}
  \leanok
  In a real inner-product space, for all $e,f$ and all real $Z$,
  \[
    \norm{e-Zf}^2=\norm{e+Zf}^2-4Z\,\langle f,e\rangle .
  \]
\end{lemma}

\begin{proof}
  \leanok
  Expand both squares.
\end{proof}

\begin{theorem}[T8, the bound for any effort and flow]\label{thm:T8}
  \lean{Atlas.wave_bound, Atlas.cayley_bound}
  \leanok
  \uses{lem:T8-identity}
  Let $Z>0$, $c\ge0$, and let $e,f$ satisfy
  $\langle f,e\rangle\ge c\norm{e}^2$ and $\norm{f}\le M\norm{e}$.  Then
  \[
    \norm{e-Zf}^2\le\Bigl(1-\frac{4Zc}{(1+ZM)^2}\Bigr)\norm{e+Zf}^2 .
  \]
  In particular, for a bounded linear $\Lambda$ with
  $\langle\Lambda e,e\rangle\ge c\norm{e}^2$ for all $e$ and
  $\norm{\Lambda}\le M$, this holds with $f=\Lambda e$ for every $e$.
\end{theorem}

\begin{proof}
  \leanok
  \uses{lem:T8-identity}
  With $y=e+Zf$: $\norm{e-Zf}^2=\norm{y}^2-4Z\langle f,e\rangle
  \le\norm{y}^2-4Zc\norm{e}^2$, and $\norm{y}\le(1+ZM)\norm{e}$.
\end{proof}

\begin{theorem}[T8, the scattering map]\label{thm:T8-cayley}
  \lean{Atlas.cayley_transform, Atlas.cayley_nonexpansive}
  \leanok
  \uses{thm:T8}
  Let $E$ be a finite-dimensional real inner-product space, not the zero
  space, and $\Lambda$ as in Theorem~\ref{thm:T8}.  Then $I+Z\Lambda$ is
  invertible and $\mathcal S=(I-Z\Lambda)(I+Z\Lambda)^{-1}$ satisfies
  \[
    \norm{\mathcal S}^2\le1-\frac{4Zc}{(1+ZM)^2}.
  \]
  If only $\langle\Lambda e,e\rangle\ge0$ for all $e$, then $I+Z\Lambda$ is
  invertible and $\norm{\mathcal S}\le1$ (on any finite-dimensional space).
\end{theorem}

\begin{proof}
  \leanok
  \uses{thm:T8}
  $\langle(I+Z\Lambda)e,e\rangle\ge\norm{e}^2$, so $I+Z\Lambda$ is injective,
  hence invertible in finite dimensions.  Every $y$ is $(I+Z\Lambda)e$ for one
  $e$, and $\mathcal Sy=(I-Z\Lambda)e$; Theorem~\ref{thm:T8}.
\end{proof}

\begin{corollary}[T8 at tier 0]\label{cor:T8-gram}
  \lean{Atlas.cayley_matrix_nonexpansive, Atlas.Piece.gram_cayley_nonexpansive}
  \leanok
  \uses{thm:T8, thm:T25-i}
  For a positive semidefinite real matrix $L$ and $Z>0$, $1+ZL$ is invertible
  and $\mathsf S=(1-ZL)(1+ZL)^{-1}$ satisfies $|\mathsf Sa|^2\le|a|^2$ for
  every $a$.  In particular this holds for the Gram port matrix
  $\tilde\Lambda$ of any network output: a tier-0 expert also serves at
  tier 2, with no certificate.
\end{corollary}

\begin{proof}
  \leanok
  \uses{thm:T8, thm:T25-i}
  The case $c=0$, with Theorem~\ref{thm:T25-i}.
\end{proof}

\section{T9a: the contraction contract, one level}

\paragraph{In plain words.}
One sweep applies every expert's scattering map to its incoming wave, hands
each outgoing wave to the neighbouring side, and adds the sources.  If every
expert's map shrinks distances by a certified factor $\rho<1$, and the
hand-over does not lengthen anything, the whole sweep shrinks distances by
$\rho$.  Then T1 applies: one answer, geometric convergence from any start,
and a distance to the classical answer set by how far each learned map is from
the classical one \emph{at the classical answer}.

\begin{definition}[The wave sweep]\label{def:wave-sweep}
  \lean{Atlas.waveSweep}
  \leanok
  Pieces $i$ in a finite set, with wave spaces $W_i$ (normed).  The global wave
  vector $a=(a_i)_i$ is measured in the $\ell^2$ product,
  $d(a,a')^2=\sum_id(a_i,a'_i)^2$.  For maps $S_i:W_i\to W_i$, an exchange
  $\Pi$ on the product and a vector $g$,
  \[
    T(a)=\Pi\bigl((S_i(a_i))_i\bigr)+g .
  \]
\end{definition}

\begin{theorem}[T9a, the sweep contracts]\label{thm:T9a-lip}
  \lean{Atlas.waveSweep_lipschitz}
  \leanok
  \uses{def:wave-sweep}
  If $d(S_ix,S_iy)\le\rho\,d(x,y)$ for every $i$, $x$, $y$, with $\rho\ge0$,
  and $d(\Pi x,\Pi y)\le d(x,y)$, then $d(Ta,Ta')\le\rho\,d(a,a')$.
\end{theorem}

\begin{proof}
  \leanok
  $d(Ta,Ta')^2\le\sum_id(S_ia_i,S_ia'_i)^2\le\rho^2\sum_id(a_i,a'_i)^2$.
\end{proof}

\begin{theorem}[T9a, the contraction contract]\label{thm:T9a}
  \lean{Atlas.contraction_contract}
  \leanok
  \uses{thm:T9a-lip, thm:T1}
  Let the $W_i$ be complete, $S_i$ the learned maps and $S^c_i$ the classical
  ones, with the same $\Pi$ and $g$, and let $a_c$ be a fixed point of the
  classical sweep.  If the hypotheses of Theorem~\ref{thm:T9a-lip} hold with
  $\rho<1$, and $d\bigl(S_i(a_{c,i}),S^c_i(a_{c,i})\bigr)\le\delta_i$ for every
  $i$, then the learned sweep has a unique fixed point $a_l$,
  $d(T^ka_0,a_l)\le\rho^kd(a_0,a_l)$ for every $a_0$ and $k$, and
  \[
    d(a_l,a_c)\le\frac{\bigl(\sum_i\delta_i^2\bigr)^{1/2}}{1-\rho}.
  \]
\end{theorem}

\begin{proof}
  \leanok
  \uses{thm:T9a-lip, thm:T1}
  Theorem~\ref{thm:T1} with the contraction of Theorem~\ref{thm:T9a-lip}; the
  defect at $a_c$ is at most $\bigl(\sum_i\delta_i^2\bigr)^{1/2}$ because
  $\Pi$ does not lengthen distances.
\end{proof}

\section{T9b: two levels}

\paragraph{In plain words.}
With a coarse space, a round is a sweep followed by a coarse solve: a
correction inside the coarse space that leaves no coarse residual.  The first
version of the design asked that each learned map contract on the
\emph{fine} space only, because the coarse space is solved exactly.  That is
not enough (Example~\ref{ex:T9b-diverges}).  What has to be certified is a
contraction factor of the \emph{composite} round.  Given that, the conclusion
is T1's, and the answer is a fixed point of the plain sweep: the coarse solve
changes how the answer is reached, not what it is.

\begin{lemma}[T9b, consistency of the coarse solve]\label{lem:T9b-consistency}
  \lean{Atlas.twoLevel_fixedPt_isFixedPt, Atlas.fixedPt_isFixedPt_twoLevel}
  \leanok
  Let $T$ be the sweep, $C$ the coarse solve and $P_0$ a linear map.  If
  $P_0(Cb-b)=Cb-b$ and $P_0\bigl(Cb-T(Cb)\bigr)=0$ for every $b$, then
  $C(Ta)=a$ implies $Ta=a$.  Conversely, if $Cb=b$ whenever
  $P_0(b-Tb)=0$, then $Ta=a$ implies $C(Ta)=a$.
\end{lemma}

\begin{proof}
  \leanok
  Let $b=Ta$ and $a=Cb$.  Then $a-b$ is fixed by $P_0$, and
  $P_0(a-Ta)=0$; but $a-Ta=a-b$, so $a-b=0$.  The converse is immediate.
\end{proof}

\begin{theorem}[T9b, the two-level contract]\label{thm:T9b}
  \lean{Atlas.two_level_contract}
  \leanok
  \uses{lem:T9b-consistency, thm:T1}
  On a complete normed space, let $T_l,C_l$ be the learned sweep and coarse
  solve, satisfying the hypotheses of Lemma~\ref{lem:T9b-consistency}, and
  $T_c,C_c$ the classical ones, with $C_c(T_ca_c)=a_c$.  If
  $d\bigl(C_lT_lx,C_lT_ly\bigr)\le\rho\,d(x,y)$ for all $x,y$ with
  $0\le\rho<1$, and $d\bigl(C_lT_la_c,C_cT_ca_c\bigr)\le\delta$, then there is
  a unique $a_l$ with $C_l(T_la_l)=a_l$; it satisfies $T_la_l=a_l$;
  $d\bigl((C_lT_l)^ka_0,a_l\bigr)\le\rho^kd(a_0,a_l)$; and
  $d(a_l,a_c)\le\delta/(1-\rho)$.
\end{theorem}

\begin{proof}
  \leanok
  \uses{lem:T9b-consistency, thm:T1}
  Theorem~\ref{thm:T1} for the composite maps, then
  Lemma~\ref{lem:T9b-consistency}.
\end{proof}

\begin{example}[Contracting the fine space is not enough]\label{ex:T9b-diverges}
  \lean{Atlas.two_level_fine_contraction_diverges}
  \leanok
  On $\R^2$ with coarse space the first axis and fine space the second, let
  \[
    G=\begin{pmatrix}0&\tfrac12\\10&0\end{pmatrix},\qquad
    C(x,y)=\bigl(\tfrac12y,\ y\bigr).
  \]
  Then $G$ contracts the fine space by $\tfrac12$; $C$ changes only the coarse
  component and leaves no coarse residual, $\bigl(Ce-G(Ce)\bigr)_1=0$; the only
  fixed point of $G$ is $0$; and
  $(CG)^k(1,2)=\bigl(5^k,\ 2\cdot5^k\bigr)$ for every $k$.
\end{example}

\begin{proof}
  \leanok
  $CG(x,y)=(5x,10x)$.
\end{proof}

\section{A remark beside T1: accuracy is not a substitute}

\paragraph{In plain words.}
T1 asks that the learned map contract.  It cannot be replaced by ``the
classical map contracts and the learned map is uniformly close to it''.  The
example below satisfies both of those, and its learned iteration alternates
between two points for ever.

\begin{example}[The SNI counterexample]\label{ex:sni}
  \lean{Atlas.sni_counterexample}
  \leanok
  \uses{thm:T1}
  Let $T(u)=\tfrac12u$ and $\tilde T(u)=\tfrac12\bigl(u-\tanh10u\bigr)$ on $\R$.
  Then $T$ is a contraction with factor $\tfrac12$ and fixed point $0$;
  $|\tilde T(u)-T(u)|\le\tfrac12$ for every $u$; and there is an
  $x\in[0.3,\,0.34]$ with $\tilde T(x)=-x$ and $\tilde T(-x)=x$, so that
  $\tilde T^k(x)=(-1)^kx$ for every $k$.
\end{example}

\begin{proof}
  \leanok
  $\tilde T(x)=-x$ means $\tanh10x=3x$.  The function $\tanh10x-3x$ is
  continuous, positive at $0.3$ (since $e^6>19$ gives $\tanh3>0.9$) and
  negative at $0.34$ (since $\tanh<1<1.02$); the intermediate value theorem
  gives $x$.  $\tilde T$ is odd.
\end{proof}
